\clearpage

\section{Code Property Graph(CPG) 자동 생성 기술}\label{sec:cpg}

본 문서 \ref{sec:cpg}절에서는 오픈소스 Joern 기반의 소스코드 그래프화에 대하여 기술한다.

\subsection{Joern CPG}

\subsubsection{Joern}
소스코드는 대체로 자연어 기반의 개발 친화적으로 비정형 데이터로 분류된다.
이러한 비정형 자연어는 정형화된 데이터를 처리하는 AI에게 가장 큰 병목으로, 이를 규격화하기 위한 시도들이 계속되고 있다.
이러한 접근들 중 하나로 오픈소스인 Joern이 존재한다.
Joern은 Java 와 각 언어의 컴파일러 기반의 명령어 프로그램으로 소스코드를 다양한 형태의 \textbf{\abbrCPG}로 제공한다.
Joern을 사용함으로써 가장 큰 장벽으로 Joern이 제공하는 그래프의 속성과 형태 그리고 다양한 Node와 Edge로 인해 분석에 어려움이 있다.
Joern에서 제공하는 다양한 Node, Edge, Properties 중에 본 기술에서 사용된 일부는 \ref{subsubsec:cpgspec}에 기술 되어있다.
이러한 복잡한 구조를 해결하기 위해 조금 더 사용자 친화적이고 AI가 처리할 수 있는 방향으로 전처리를 하게된다.
이렇게 전처리 후에 생성된 그래프를 본 문서에서는 \textbf{Template(템플릿)}이라고 규정한다.

\subsubsection{CPG Specification}\label{subsubsec:cpgspec}
CPG 명세 1.1은 정적 분석을 위해 설계된 언어 중립적인 intermediate graph representation을 정의한다. CPG는 방향성 directed, 엣지 라벨 edge-labeled, 속성 attributed 멀티그래프 구조를 가지며, 코드의 각기 다른 측면을 모델링하는 여러 개의 레이어 Layer로 구성된다.

\paragraph{MetaData 레이어} 이 레이어는 그래프의 생성 정보와 출처(provenance)를 기술한다.
\begin{itemize} \setlength{\itemsep}{0pt} \setlength{\parskip}{0pt} \setlength{\topsep}{0pt} \setlength{\partopsep}{0pt}
    \item \textbf{노드:} \texttt{META\_DATA} (언어, 명세 버전, 오버레이 정보 포함).
    \item \textbf{속성:} \texttt{LANGUAGE}, \texttt{VERSION}, \texttt{OVERLAYS}, \texttt{ROOT}, \texttt{HASH}.
\end{itemize}

\paragraph{FileSystem 레이어} 그래프 노드를 실제 물리적 파일 시스템의 위치와 매핑한다.
\begin{itemize} \setlength{\itemsep}{0pt} \setlength{\parskip}{0pt} \setlength{\topsep}{0pt} \setlength{\partopsep}{0pt}
    \item \textbf{노드:} \texttt{FILE} (소스 파일을 표현).
    \item \textbf{엣지:} \texttt{SOURCE\_FILE} (노드를 원본 파일과 연결).
    \item \textbf{속성:} \texttt{FILENAME}, \texttt{CODE}, \texttt{COLUMN\_NUMBER}, \texttt{LINE\_NUMBER}.
\end{itemize}

\paragraph{Namespace 레이어} 패키징 및 논리적 네임스페이스 구조를 다룬다.
\begin{itemize} \setlength{\itemsep}{0pt} \setlength{\parskip}{0pt} \setlength{\topsep}{0pt} \setlength{\partopsep}{0pt}
    \item \textbf{노드:} \texttt{NAMESPACE} (논리적 네임스페이스), \texttt{NAMESPACE\_BLOCK} (네임스페이스 내의 구체적인 코드 블록).
    \item \textbf{속성:} \texttt{NAME}, \texttt{FULL\_NAME}.
\end{itemize}

\paragraph{Method 레이어} 파라미터와 반환 타입을 포함하여 메서드 선언을 모델링한다.
\begin{itemize} \setlength{\itemsep}{0pt} \setlength{\parskip}{0pt} \setlength{\topsep}{0pt} \setlength{\partopsep}{0pt}
    \item \textbf{노드:} \texttt{METHOD} (함수 자체), \texttt{METHOD\_PARAMETER\_IN}, \texttt{METHOD\_PARAMETER\_OUT}, \texttt{METHOD\_RETURN}.
    \item \textbf{속성:} \texttt{SIGNATURE}, \texttt{IS\_EXTERNAL}, \texttt{FULL\_NAME}, \texttt{IS\_VARIADIC}.
\end{itemize}

\paragraph{Type 레이어} 타입의 계층 구조, 선언 및 사용을 표현한다.
\begin{itemize} \setlength{\itemsep}{0pt} \setlength{\parskip}{0pt} \setlength{\topsep}{0pt} \setlength{\partopsep}{0pt}
    \item \textbf{노드:} \texttt{TYPE} (구체적 인스턴스), \texttt{TYPE\_DECL} (클래스/인터페이스 선언), \texttt{MEMBER} (필드), \texttt{TYPE\_ARGUMENT}.
    \item \textbf{엣지:} \texttt{BINDS\_TO}, \texttt{INHERITS\_FROM}.
    \item \textbf{속성:} \texttt{TYPE\_FULL\_NAME}, \texttt{INHERITS\_FROM\_TYPE\_FULL\_NAME}.
\end{itemize}

\paragraph{\abbrAST 레이어} 코드의 syntactic structure를 표현한다.
\begin{itemize} \setlength{\itemsep}{0pt} \setlength{\parskip}{0pt} \setlength{\topsep}{0pt} \setlength{\partopsep}{0pt}
    \item \textbf{노드:} \texttt{BLOCK}, \texttt{CALL}, \texttt{CONTROL\_STRUCTURE}, \texttt{IDENTIFIER}, \texttt{LITERAL}, \texttt{RETURN}, \texttt{UNKNOWN}, \texttt{MODIFIER}.
    \item \textbf{엣지:} \texttt{AST} (부모 노드와 자식 노드를 연결).
    \item \textbf{속성:} \texttt{ORDER}, \texttt{CODE}.
\end{itemize}

\paragraph{CallGraph 레이어} 메서드 호출 및 dispatching 과정을 모델링한다.
\begin{itemize} \setlength{\itemsep}{0pt} \setlength{\parskip}{0pt} \setlength{\topsep}{0pt} \setlength{\partopsep}{0pt}
    \item \textbf{노드:} \texttt{CALL} (호출 지점).
    \item \textbf{엣지:} \texttt{CALL} (발신 호출), \texttt{RECEIVER}, \texttt{ARGUMENT}.
    \item \textbf{속성:} \texttt{METHOD\_FULL\_NAME}, \texttt{DISPATCH\_TYPE}, \texttt{ARGUMENT\_INDEX}.
\end{itemize}

\paragraph{\abbrCFG 레이어} 메서드 내에서의 실행 흐름(execution flow)을 모델링한다.
\begin{itemize} \setlength{\itemsep}{0pt} \setlength{\parskip}{0pt} \setlength{\topsep}{0pt} \setlength{\partopsep}{0pt}
    \item \textbf{노드:} \texttt{CFG\_NODE} (모든 제어 흐름 노드의 기본 클래스).
    \item \textbf{엣지:} \texttt{CFG} (실행 순서를 지시).
\end{itemize}

\paragraph{\abbrPDG 레이어} 오염 분석 taint analysis 및 슬라이싱을 위해 데이터 및 제어 의존성을 모델링한다.
\begin{itemize} \setlength{\itemsep}{0pt} \setlength{\parskip}{0pt} \setlength{\topsep}{0pt} \setlength{\partopsep}{0pt}
    \item \textbf{엣지:} 
    \begin{itemize} \setlength{\itemsep}{0pt} \setlength{\parskip}{0pt} \setlength{\topsep}{0pt} \setlength{\partopsep}{0pt}
        \item \texttt{CDG} (제어 의존성 그래프, Control Dependence Graph).
        \item \texttt{REACHING\_DEF} (데이터 의존성).
    \end{itemize}
    \item \textbf{속성:} \texttt{VARIABLE} (전파되는 변수).
\end{itemize}

\paragraph{Binding 레이어} 메서드 이름 및 시그니처를 타입 선언과 resolve한다.
\begin{itemize} \setlength{\itemsep}{0pt} \setlength{\parskip}{0pt} \setlength{\topsep}{0pt} \setlength{\partopsep}{0pt}
    \item \textbf{노드:} \texttt{BINDING} (이름/시그니처를 메서드와 연결).
    \item \textbf{엣지:} \texttt{BINDS} (TypeDecl에서 Binding으로), \texttt{REF} (Binding에서 Method로).
\end{itemize}

\paragraph{보조 Auxiliary Layers} 특정 분석 요구사항을 충족하기 위한 추가 레이어들이 존재한다.
\begin{itemize} \setlength{\itemsep}{0pt} \setlength{\parskip}{0pt} \setlength{\topsep}{0pt} \setlength{\partopsep}{0pt}
    \item \textbf{Dominators:} \texttt{DOMINATE} 및 \texttt{POST\_DOMINATE} 엣지를 포함한다.
    \item \textbf{Finding:} 분석 결과 저장을 위한 노드 (\texttt{FINDING}, \texttt{KEY\_VALUE\_PAIR}).
    \item \textbf{TagsAndLocation:} 임의의 노드 태깅을 위한 \texttt{TAG} 및 \texttt{LOCATION} 노드.
\end{itemize}

각 레이어별 상세 스키마, 노드 속성, 그리고 엣지 제약 조건에 대한 보다 구체적인 정보는 CPG 공식 웹사이트(\url{https://cpg.joern.io/})에서 확인할 수 있다.

\subsection{가상 컨테이너}

Joern은 활발히 개발·배포되는 도구로, 버전 변화에 따라 생성되는 \abbrCPG 스키마와 속성이 달라질 수 있다.
분석 파이프라인의 안정성과 재현성을 확보하기 위해, 본 연구에서는 Docker 기반 가상 컨테이너를 구축하여 Joern 실행 환경(버전, 의존 패키지, 설정 값 등)을 고정하였다.
이를 통해 Template 변환 및 후속 AST/그래프 변환 모듈이 항상 동일한 형태의 CPG를 입력으로 받도록 보장하며, 장기간에 걸친 반복 실험과 대규모 배치 전처리에서도 일관된 결과를 얻을 수 있다.
도커 컨테이너 설정은 다음 Figure~\ref{fig:docker}와 같다.

\begin{figure}
    \centering
    \begin{lstlisting}[language=Python]
FROM ubuntu:22.04

ARG JOERN_VERSION=4.0.361
ENV DEBIAN_FRONTEND=noninteractive
ENV LANG=C.UTF-8 LC_ALL=C.UTF-8
ENV JAVA_HOME=/usr/lib/jvm/java-17-openjdk-amd64
ENV PATH="/usr/local/bin:${PATH}"
ENV JOERN_PORT=8080

RUN apt-get update && \
    apt-get install -y --no-install-recommends \
      openjdk-17-jdk git unzip wget curl ca-certificates \
      xauth x11-apps locales tini && \
    rm -rf /var/lib/apt/lists/* && \
    update-alternatives --install /usr/bin/java java /usr/lib/jvm/java-17-openjdk-amd64/bin/java 1

RUN set -eux; \
    mkdir -p /opt/joern; \
    wget --progress=dot:giga -O /tmp/joern-cli.zip \
      "https://github.com/joernio/joern/releases/download/v${JOERN_VERSION}/joern-cli.zip"; \
    unzip -q /tmp/joern-cli.zip -d /opt/joern; \
    rm -f /tmp/joern-cli.zip; \
    ln -sf /opt/joern/joern-cli/joern /usr/local/bin/joern; \
    test -x /usr/local/bin/joern

RUN mkdir -p /workspace && chmod 777 /workspace
WORKDIR /workspace

RUN joern --help >/dev/null 2>&1 || true; 

ENTRYPOINT ["/usr/bin/tini", "--"]
    \end{lstlisting}
    \caption{Docker 컨테이너 생성 스크립트}
    \label{fig:docker}
\end{figure}

해당 스크립트를 이용하여 도커 컨테이너를 생성한 후, 호스트 디렉토리를 컨테이너 내 \textit{/workspace}와 같은 경로에 볼륨으로 마운트하여 동작시킨다.
이와 같은 구성은 CPG 생성 절차에서 호스트와 컨테이너 간의 파일을 안정적으로 공유하게 하고, 도커에 설치된 Joern을 통해 대규모 소스 코드를 일관된 환경에서 자동 변환할 수 있도록 지원한다.

\subsection{병렬 처리 및 자동화}

하나의 소스 코드를 JSON 형태의 CPG로 생성하기 위해서는 Joern이 제공하는 두 가지 명령어를 순차적으로 실행해야 한다. \verb|joern-parse|와 \verb|joern-export|를 사용하여, 소스 코드를 컴파일과 동시에 바이너리 형태의 CPG로 변환하고 이를 JSON으로 익스포트하는 데에 각각 약 3\textasciitilde5초가 소요된다.
대량의 소스 코드를 전처리하는 경우, 직렬 방식으로는 약 1만 개의 소스 코드를 단순 CPG로 변환하는 데에만 8\textasciitilde13시간이 필요하다.
본 연구에서는 이러한 병목을 해소하기 위해 Python 기반 병렬 처리 로직을 도입하여, Joern 호출을 워커 단위로 분산 실행하였다.
그 결과, 약 21{,}000개의 데이터셋을 변환하는 데 필요한 전체 소요 시간을 8시간 수준으로 줄여, 파일당 평균 변환 시간을 \textbf{4.32초에서 1.37초로 단축}하고, 전체 처리 시간 기준으로 \textbf{68.3\%} 수준의 성능 향상을 달성하였다.


\begin{table}[H]
\centering
\begin{tabular*}{0.8\textwidth}{@{\extracolsep{\fill}} lrrrr}
\toprule
 & \makecell[r]{변환파일 수 \\ (개)} & \makecell[r]{소요 시간 \\ (분)} & \makecell[r]{파일당 소요시간 \\ (초)} & \makecell[r]{단축율 \\ (\%)} \\ \midrule
일반 & 10000 & 720 & 4.32 & \multirow{2}{*}{\textbf{-68.3\%}} \\
병렬 & 21000 & 480 & 1.37 &  \\ \bottomrule
\end{tabular*}
\caption{CPG 병렬 처리 속도 비교}
\label{tab:cpg-convertion-time-table}
\end{table}